\documentclass[12pt]{article}
\usepackage{cwilliams-standard}

\setclass{\NLP}
\settitle{Homework 1}

\begin{document}

\maketitlepage

In this report, I will go over my steps in tackling the first homework on NLP. These tasks included: 
Text-Embedding and Topic Modeling. It was also unstated (but clearly) necessary that we needed to do 
preprocessing on the corpora given.

For the first (zeroth, $0^{th}$) task, I did preprocessing. For each subsequent task, I attempted to do what I 
believed was appropriate preprocessing for said task (as we talked at length in class about how 
preprocessing was approached differently depending on what the problem you are dealing with is). I will 
do my best to explain each of the methods of preprocessing I had considered, and which ones I decided on 
using for each task and why.

\begin{aside}{Methodology}
For the two main tasks, I will explain my methodology as we go. 
\end{aside}

\begin{important}{Dataset}
    For the \textit{entire} report/assignment, I did not split the dataset up into its constituent parts 
    (train, validation, test) but, instead, \textbf{combined} them (as the assignment said to do, Pg. 1). Then, after combination 
    they are used in 5-fold cross validation. Stated for clarity.
\end{important}


\section*{Preprocessing}

Firstly, for preprocessing I laid out these options of preprocessing the data to give myself as many options 
of cleaning the data as possible.

\begin{table}[H]
    \centering
    \definecolor{tableheader}{rgb}{0.18, 0.32, 0.55}
    \definecolor{rowalt}{rgb}{0.93, 0.95, 0.98}
    \renewcommand{\arraystretch}{1.4}
    \begin{tabular}{>{\bfseries\centering\arraybackslash}p{0.6cm}
                    p{5.2cm}
                    p{7.8cm}}
        \toprule
        \rowcolor{tableheader}
        \color{white}\# & \color{white}\textbf{Step} & \color{white}\textbf{Description} \\
        \midrule
        \rowcolor{rowalt} 0 & Expand Contractions      & \texttt{"can't"} $\to$ \texttt{"cannot"} \\
                          1 & Lowercase                & Normalize all text to lowercase \\
        \rowcolor{rowalt} 2 & Remove URLs \& Mentions  & Strip \texttt{http://...} and \texttt{@username} tokens \\
                          3 & Strip Hashtag Symbol     & \texttt{\#word} $\to$ \texttt{word} (keep the term) \\
        \rowcolor{rowalt} 4 & Remove Emoji             & Drop all emoji characters \\
                          5 & Remove Punctuation       & Strip non-alphanumeric characters \\
        \rowcolor{rowalt} 6 & Remove Digit-only Tokens & Drop tokens consisting entirely of digits \\
                          7 & Remove Stop Words        & Filter standard English stop words \\
        \rowcolor{rowalt} 8 & POS-aware Lemmatization  & WordNet lemmatization guided by POS tags \\
        \bottomrule
    \end{tabular}
    \caption{Preprocessing Options}
    \label{tab:preprocessing}
\end{table}

These options were taken from class with one exception: \# 8, POS-aware Lemmatization. 
This was suggested by Claude Code (within VSCode) to tag parts of speech in the preprocesing stage. 
This specific lemmatization is aware of parts of speech (i.e., understands that "running" 
is a verb and will convert that to its root, "run") but does not pass along any additional 
information about the word to the frequency-based embedder.

Each preprocessing option that was used in each task will be stated.

\section*{Text-Embedding}

\section{Frequency-based Embeddings}

I used \textit{both} Bag of Words and TF-IDF to compare the two methods and see which resulted 
in better performance.

\section{Word-based Embeddings}

\section{Sentence Embeddings}

\section*{Topic Modeling}

\section*{Findings}

\end{document}